%% =====================================================================
\section{Background}
\label{sec:background}
%% =====================================================================

This section introduces the concepts on which the rest of the paper depends: caching
patterns and expiry semantics, the cache stampede, overload behavior of connection
pools, rate-limiting algorithms, unsupervised anomaly detection with isolation
forests, and the container and observability tooling used by the testbed.

\subsection{Caching Patterns and Expiry Semantics}
\label{sec:bg-caching}
Three caching patterns dominate web back ends \cite{kleppmann2017}:
\begin{itemize}[leftmargin=*]
\item \textbf{Cache-aside (lazy loading).} The application reads the cache; on a miss
it reads the database and writes the result into the cache. The cache never talks to
the database itself. This is the pattern used by CacheCollapse and by the memcache
deployment described by Nishtala et al.\ \cite{nishtala2013}.
\item \textbf{Read-through.} The cache library itself loads missing values from the
database, which centralizes regeneration but does not by itself prevent concurrent
regeneration of the same key.
\item \textbf{Write-through / write-behind.} Writes update the cache and the database
together (or asynchronously), keeping entries fresh and reducing reliance on expiry.
\end{itemize}
Expiry is usually governed by a TTL. Short TTLs bound staleness but increase the
rate of expiries; long TTLs reduce database load but serve older data. HTTP caching
formalizes a useful middle ground: the \texttt{stale-while-revalidate} directive
allows a cache to serve a stale response for a bounded period while it refreshes the
entry in the background \cite{rfc5861}. CacheCollapse applies the same idea inside
the application tier.

Popularity in real workloads is strongly skewed; web object requests follow a
Zipf-like distribution \cite{breslau1999}, and large key--value workloads concentrate
a large fraction of requests on a small set of keys \cite{atikoglu2012}. Skew is what
makes caching effective, and also what makes a few hot keys a single point of
failure: Fan et al.\ show that a small front-end cache is needed precisely to prevent
hot keys from overloading individual back-end nodes \cite{fan2011}.

\subsection{The Cache Stampede}
\label{sec:bg-stampede}
Consider a hot key $k$ with TTL $T$ and regeneration time $d$ (the time to compute the
value from the database). Let requests for $k$ arrive at rate $\lambda_k$. If the
entry expires at time $t_{\mathrm{exp}}$, every request arriving in
$(t_{\mathrm{exp}}, t_{\mathrm{exp}}+d)$ finds the cache empty. In a naive cache-aside
implementation each of those requests independently regenerates the value, so the
expected number of database queries caused by one expiry is
\begin{equation}
\mathbb{E}[Q_{\mathrm{exp}}] \approx 1 + \lambda_k d ,
\label{eq:stampede}
\end{equation}
and for a synchronized burst of $B$ concurrent requests it is simply $B$. Equation
\eqref{eq:stampede} has two important properties. First, the redundant work grows
with $d$, and $d$ itself grows as the database becomes loaded---a positive feedback
loop. Second, it is \emph{independent of the TTL}; a longer TTL makes stampedes less
frequent but not smaller.

Four families of mitigation exist:
(i)~\emph{request coalescing} (single-flight), in which only one caller regenerates a
key and the others wait for or reuse its result \cite{nishtala2013,singleflight};
(ii)~\emph{serving stale data} during regeneration \cite{rfc5861,nygard2018};
(iii)~\emph{early or probabilistic recomputation} before expiry, which Vattani et al.\
show can be made optimal with the XFetch rule \cite{vattani2015}; and
(iv)~\emph{TTL jitter}, which de-synchronizes expiries of keys that were filled at the
same time. Only the first two remove redundancy once a synchronized burst has already
arrived; the latter two reduce the probability that it happens.

\subsection{Connection Pools, Queueing, and Overload}
\label{sec:bg-queueing}
An application reaches its database through a bounded pool of $c$ connections. If each
query holds a connection for $d$ seconds, the pool's nominal service capacity is
\begin{equation}
\mu_{\max} = \frac{c}{d}\quad\text{queries/s.}
\label{eq:capacity}
\end{equation}
Little's law, $L=\lambda W$ \cite{little1961}, relates the mean number of in-flight
requests $L$ to the arrival rate $\lambda$ and mean residence time $W$; as $\lambda$
approaches $\mu_{\max}$ the residence time---and thus latency---grows without bound in
classical queueing models \cite{kleinrock1975}. Pools differ in what happens at
saturation. A \emph{blocking} pool queues callers, converting overload into latency;
a \emph{fail-fast} pool rejects them, converting overload into errors. The psycopg2
\texttt{ThreadedConnectionPool} used here is fail-fast: requesting a connection from
an exhausted pool raises an exception immediately \cite{psycopgpool}.

A well-known pathology of overloaded systems is that useful throughput (\emph{goodput})
can \emph{fall} as offered load rises, because more and more work is spent on requests
that will ultimately be rejected or time out. This is the essence of congestion
collapse in networks \cite{jacobson1988} and of receive livelock in operating systems
\cite{mogul1997}. Overload-control research therefore argues for shedding excess load
\emph{early}, at the edge of the system, before it consumes scarce resources
\cite{welsh2001,welsh2003,cherkasova2002,beyer2016}. Tail latency compounds the
problem: in fan-out services, a small fraction of slow components dominates
user-perceived latency \cite{dean2013}.

\subsection{Rate-Limiting Algorithms}
\label{sec:bg-ratelimit}
Rate limiting caps the rate at which a client (or a class of clients) may consume a
resource. Table~\ref{tab:ratelimit-algos} summarizes the algorithms relevant to this
work. The \emph{leaky bucket} smooths output to a constant rate and is the basis of
Nginx's \texttt{limit\_req} module \cite{nginxlimitreq,tanenbaum2011}; a configurable
burst allowance lets short spikes through without delay (\texttt{nodelay}). A
\emph{fixed-window counter} is the simplest to implement on a shared store such as
Redis: a counter keyed by client and time window is incremented on each request and
compared with a threshold. Rate limiting can be enforced locally or across a cluster
\cite{raghavan2007}. HTTP defines status 429 (Too Many Requests) for requests rejected
by a rate limiter \cite{rfc6585}, distinguishing intentional shedding from server
failure (status 503).

\begin{table}[t]
\centering
\caption{Rate-limiting algorithms relevant to CacheCollapse.}
\label{tab:ratelimit-algos}
\small
\begin{tabularx}{\linewidth}{@{}L{2.6cm}YYL{2.6cm}@{}}
\toprule
Algorithm & Mechanism & Strengths / weaknesses & Used here \\
\midrule
Fixed window & Counter per client per window; reject above limit & Trivial on Redis; allows up to $2\times$ limit across a window boundary & Flask app layer (1~s window, 50 req/s) \\
Sliding window & Weighted counts over the last window & Smoother than fixed window; more state & -- \\
Token bucket & Tokens accrue at rate $r$ up to depth $b$; a request spends one & Permits bursts up to $b$; average rate $r$ & -- \\
Leaky bucket & Requests drain from a queue at constant rate $r$ & Smooth output; with \texttt{burst}/\texttt{nodelay} behaves like a token bucket & Nginx edge (\texttt{limit\_req}, 50 r/s, burst 20) \\
Admission control & Admit/reject by measured load & Adapts to capacity; needs a load signal & Discussed (Sec.~\ref{sec:discussion}) \\
\bottomrule
\end{tabularx}
\end{table}

\subsection{Anomaly Detection and the Isolation Forest}
\label{sec:bg-iforest}
Intrusion detection has long been split between \emph{misuse} (signature or rule)
detection and \emph{anomaly} detection, which models normal behavior and flags
deviations \cite{denning1987}. Anomaly detection is attractive for attacks without a
fixed signature, but it is notoriously sensitive to the base rate of attacks
\cite{axelsson2000}, and machine-learning detectors that perform well in the
laboratory often fail operationally because ``anomalous'' is not the same as
``malicious'' \cite{sommer2010}. Surveys organize outlier detection by technique and
by the availability of labels \cite{chandola2009,hodge2004}.

The \emph{Isolation Forest} (iForest) \cite{liu2008,liu2012} is an unsupervised
ensemble method that isolates anomalies directly rather than profiling normal points.
Each of $t$ trees is grown on a sub-sample of $\psi$ points by repeatedly choosing a
random feature and a random split value between that feature's minimum and maximum,
until each point is isolated or a height limit is reached. Anomalies, being few and
different, are isolated near the root, so their path length $h(x)$ is short. The
anomaly score of a point $x$ is
\begin{equation}
s(x,\psi) = 2^{-\frac{\mathbb{E}[h(x)]}{c(\psi)}}, \qquad
c(\psi) = 2H(\psi-1) - \frac{2(\psi-1)}{\psi},
\label{eq:iforest}
\end{equation}
where $\mathbb{E}[h(x)]$ is the mean path length over the trees, $H(i)\approx\ln i +
0.5772$ is the harmonic number, and $c(\psi)$ normalizes by the average path length of
an unsuccessful binary-search-tree lookup. Scores near 1 indicate anomalies; scores
well below 0.5 indicate normal points. The scikit-learn implementation
\cite{pedregosa2011} exposes a \texttt{decision\_function} (positive for inliers,
negative for outliers) whose threshold is set by a \texttt{contamination} parameter,
the expected fraction of outliers in the training data. Isolation forests require no
labels, train in near-linear time, and work well with small samples, which makes them
suitable for learning a ``normal'' profile from a few seconds of live traffic. Unlike
the supervised random forests from which they borrow the ensemble idea
\cite{breiman2001}, they cannot learn what an attack looks like---only what normal
looks like.

\subsection{Containers and Observability}
\label{sec:bg-containers}
Containers package an application with its dependencies and isolate it using kernel
namespaces and control groups \cite{merkel2014,pahl2015}; orchestrators from Docker
Compose \cite{dockercompose} to Kubernetes \cite{burns2016} wire many containers into
a microservice application \cite{dragoni2017}. Containers bring their own security
concerns, chiefly around isolation and configuration \cite{combe2016,sultan2019}, but
for this study their key property is \emph{reproducibility}: an entire multi-service
system can be created, attacked, measured, and destroyed from a single declarative
file. Observability is provided by Prometheus, a pull-based time-series system that
periodically scrapes text-format metrics endpoints \cite{brazil2018}, visualized with
Grafana. For low-latency browser dashboards, Server-Sent Events (SSE) provide a simple
one-way push channel over HTTP.
